% Figure: Hall-effect thruster cross-section. An annular discharge channel
% carries an axial electric field and a radial magnetic field; the crossed
% fields trap electrons in an azimuthal Hall current while xenon ions are
% accelerated out of the channel as the beam. The end-to-end case study
% carries the numbers.
\begin{figure}[htbp]
\centering
\begin{tikzpicture}[scale=1.0]
  % discharge channel walls (annular, shown in cross-section: two boxes)
  \fill[enggray!15] (0,0.7) rectangle (3.2,1.5);
  \fill[enggray!15] (0,-1.5) rectangle (3.2,-0.7);
  \draw[enggray, thick] (0,0.7) rectangle (3.2,1.5);
  \draw[enggray, thick] (0,-1.5) rectangle (3.2,-0.7);
  \node[enggray, font=\scriptsize, anchor=west] at (0.1,1.1) {dielectric wall};
  % anode at the back of the channel
  \draw[engred, line width=1.2pt] (0.15,-0.65) -- (0.15,0.65);
  \node[engred, font=\scriptsize, anchor=east, rotate=90] at (0.05,0) {anode (+)};
  % magnetic coils / poles producing radial B
  \node[engblue, font=\scriptsize, anchor=south] at (2.5,1.55) {N};
  \node[engblue, font=\scriptsize, anchor=north] at (2.5,-1.55) {S};
  % radial B field arrows (across the channel)
  \foreach \x in {1.4,2.0,2.6} {
    \draw[engblue, ->, thick] (\x,0.65) -- (\x,-0.6);
  }
  \node[engblue, font=\scriptsize, anchor=south west] at (2.6,0.0) {$\vect{B}$ (radial)};
  % axial E field arrow inside the channel
  \draw[engorange, ->, line width=1pt] (0.4,0) -- (3.0,0);
  \node[engorange, font=\scriptsize, anchor=south] at (1.6,0.05) {$\vect{E}$ (axial)};
  % azimuthal Hall electron current (into/out of page markers)
  \node[englime, font=\scriptsize] at (1.9,0.35) {$\otimes$};
  \node[englime, font=\scriptsize] at (1.9,-0.35) {$\odot$};
  \node[englime, font=\scriptsize, anchor=west] at (3.25,0.4)
    {Hall electron current};
  % accelerated ion beam leaving the channel
  \foreach \y in {-0.35,0,0.35} {
    \draw[engred, ->, thick] (3.0,\y) -- (4.3,\y);
  }
  \node[engred, font=\scriptsize, anchor=west] at (4.0,-0.7) {ion beam};
  % cathode neutraliser
  \fill[engblack] (4.2,1.1) circle (0.08);
  \node[engblack, font=\scriptsize, anchor=west] at (4.3,1.1)
    {cathode (neutraliser)};
  \draw[engblack, ->, dashed] (4.2,1.0) .. controls (4.0,0.3) .. (3.8,0.05);
\end{tikzpicture}
\caption[Hall-effect thruster]{Cross-section of one side of an annular
Hall-effect thruster. An anode at the back of the dielectric channel sets up
an axial electric field $\vect{E}$, while magnet coils drive a radial
magnetic field $\vect{B}$ across the channel. The crossed fields trap the
light electrons in an azimuthal $E\times B$ drift, the Hall current, which
greatly raises their residence time and lets them ionise the injected xenon
efficiently. The heavy ions have a Larmor radius larger than the channel and
are not magnetised, so they fall freely down the axial field and leave as a
high-velocity beam. An external cathode supplies electrons both to sustain
the discharge and to neutralise the departing beam, without which the
spacecraft would charge negative and pull the ions back. The end-to-end case
study carries the discharge, beam, thrust, and efficiency numbers.}
\label{fig:vol04ch13:hall-thruster}
\end{figure}
